\section{Algorithm dossiers: every method, ten questions deep}
\label{sec:dossiers}

The body sections answered the consolidated question \emph{which method
wins, and why}. This section answers the \emph{disaggregated} question
that an adversarial reviewer asks: \emph{for each individual method, does
the report explain it well enough that a careful student can re-derive
its behaviour on \mimic{} from first principles, find it in the code,
predict the failure mode before reading the result, and recover the same
deployment recommendation?} Each subsection is structured as ten questions:
\textbf{(Q1)} formal definition, \textbf{(Q2)} mechanism hypothesis,
\textbf{(Q3)} code location, \textbf{(Q4)} mathematical prediction,
\textbf{(Q5)} observation, \textbf{(Q6)} prediction-vs-observation
reconciliation, \textbf{(Q7)} failure mode, \textbf{(Q8)} fix or
mitigation, \textbf{(Q9)} deployment guidance, \textbf{(Q10)} what we
would do differently. The dossiers are dense by design; they are the
self-contained reading guide for every algorithmic axis of the report.

\subsection{Dossier 1 -- \fedavg{} (the floor everything must beat)}
\label{sec:dossier-fedavg}

\textbf{Q1. Formal definition.}
At round $t$, sample $S_t \subseteq \{1,\dots,K\}$ of size $C$ uniformly
without replacement. Each $k \in S_t$ receives the global state $w_t$,
runs $E$ epochs of mini-batch SGD/Adam on its private dataset
$\mathcal{D}_k$ with learning rate $\eta$ and batch size $B$, producing
$w_{t,k}^{(E)}$. The server aggregates with size weights
$p_k = |\mathcal{D}_k| / \sum_{j \in S_t} |\mathcal{D}_j|$:
\[
w_{t+1} \;=\; \sum_{k \in S_t} p_k \, w_{t,k}^{(E)}.
\]
Default knobs in this study: $C\!=\!5$, $E\!=\!1$, $B\!=\!256$,
$\eta\!=\!10^{-3}$, Adam.

\textbf{Q2. Mechanism hypothesis.}
\fedavg{} is the size-weighted consensus algorithm. With $E\!=\!1$ and
moderate $\eta$ the per-client step is small relative to the global
geometry, so the size-weighted average is approximately a
\emph{population-weighted} gradient step. We therefore predict that
\fedavg{} will track the centralized solution on dominant clients
(MICU, MICU/SICU, CVICU) at the cost of small clients (Neuro
Intermediate, Neuro Stepdown), expressed as low worst-client recall.

\textbf{Q3. Code location.} \texttt{flopt/fedavg.py}, helpers
\texttt{flopt/training\_helpers.py}; aggregation is in
\texttt{flopt/aggregators.py}. The only delta from textbook \fedavg{}
is the size-weighted (rather than uniform) aggregator -- chosen because
size weighting is the more honest baseline given the natural ICU sizes.

\textbf{Q4. Math prediction.}
With one local epoch and Adam, the local update for client $k$ at round
$t$ is approximately a single Adam step on $F_k$ from $w_t$, scaled by
$|\mathcal{D}_k|/B$ minibatches. After size-weighted averaging the
expected drift per round is
$\mathbb{E}\| w_{t+1} \!-\! w_t \|_2 \approx \eta \cdot
\bar{n}\,\|\nabla F(w_t)\|_2$ where $\bar{n} = (\sum_k p_k\,n_k)/B$.
Taking $n \!=\! |\mathcal{D}_k|/B$, the dominant clients (MICU
$\sim$$11{,}000$ rows) contribute $\sim$$43$ minibatches per round
versus $\sim$$2$ for Neuro Intermediate ($\sim$$500$ rows). Thus the
size-weighted average is dominated by MICU's gradient direction by
roughly $20{:}1$, i.e.\ \fedavg{} should look like ``MICU centralized.''

\textbf{Q5. Observation.}
Headline: $\auprc{} \!=\! 0.632 \!\pm\! 0.014$,
$\worstrec{} \!=\! 0.21 \!\pm\! 0.14$ (Section~\ref{sec:results},
Table~\ref{tab:headline}). Per-client clinical metrics
(\texttt{outputs/full\_mimic\_iv\_training/metrics/per\_client\_clinical\_metrics.csv})
show MICU sensitivity $0.835$, Neuro Intermediate sensitivity $0.40$.
The dominant clients win, the small clients lose -- exactly the
prediction.

\textbf{Q6. Reconciliation.}
The $20{:}1$ minibatch ratio computed in Q4 predicts that gradients on
MICU will dominate every step. The observed Neuro Intermediate
sensitivity of $0.4$ on $10$ deaths out of $504$ rows is the predicted
failure mode in calibration form: the model does not see enough Neuro
deaths to push its decision boundary toward the rare class for that
unit. Drift telemetry
(\texttt{outputs/full\_mimic\_iv\_training/metrics/method\_summary.csv})
confirms the per-round drift is $0.013$--$0.015$ -- small per round,
but persistently MICU-pulled.

\textbf{Q7. Failure mode.}
\fedavg{} fails on \emph{worst-client recall} and on the small Neuro
units specifically. The mode is \emph{not} optimization failure (loss
converges, gradients are stable); it is \emph{population-weighted
representation}. The model is \emph{accurate on average}, which is the
wrong objective for a fairness-constrained ICU deployment.

\textbf{Q8. What we changed to fix it.}
We did \emph{not} modify \fedavg{}. Instead, we layered \fedprox{} (Q21
below) and \cvar{} on top, and added a \emph{logistic} backbone as a
calibration sanity check. \fedprox{}-$\mu\!=\!0.1$ raises worst-client
recall to $0.50 \!\pm\! 0.05$, more than doubling \fedavg{}'s $0.21$
without sacrificing average performance.

\textbf{Q9. Deployment guidance.}
Use \fedavg{} only when (a) clients are roughly homogeneous in label
distribution and size, or (b) average accuracy is the stated
objective. On a heterogeneous ICU like \mimic{}, \fedavg{}
\emph{operationally} writes off the small Neuro units; do not deploy
\fedavg{} alone if any client constituency is small.

\textbf{Q10. What we would do differently.}
Move from one-epoch local training to adaptive $E$ per client (small
clients get more local passes), use loss-weighted (rather than
size-weighted) aggregation, and measure per-client \auprc{} as the
primary objective rather than headline accuracy.

\subsection{Dossier 2 -- \fedprox{} (the proximal anchor)}
\label{sec:dossier-fedprox}

\textbf{Q1. Formal definition.}
Each client minimises
$F_k^{\mathrm{prox}}(w; w_t) = F_k(w) + (\mu/2)\, \|w - w_t\|_2^2$ for
$E$ local steps. Aggregation is unchanged from \fedavg{}.

\textbf{Q2. Mechanism hypothesis.}
The proximal term is a \emph{quadratic anchor} that prevents the local
solution from straying too far from $w_t$. On heterogeneous data,
unconstrained local steps drift toward each client's own minimum
$w_k^*$; the proximal term costs the trip in the local objective
itself, so the size-weighted average is a less noisy consensus.

\textbf{Q3. Code location.} \texttt{flopt/fedprox.py}, function
\texttt{train\_one\_client\_fedprox}. The proximal term iterates over
\texttt{model.named\_parameters()} (excluding BatchNorm/dropout
buffers) -- a correctness fix from an earlier version that included
buffers and produced silently wrong $\mu$ effects.

\textbf{Q4. Math prediction.}
Adding $(\mu/2)\|w - w_t\|^2$ to the local loss makes each minibatch
step
$w \!\leftarrow\! w - \eta(\nabla F_k(w) + \mu(w - w_t))$.
At equilibrium of one local pass,
$w^{(E)} \!\approx\! w_t - \eta_{\mathrm{eff}} \nabla F_k(w_t)$,
where $\eta_{\mathrm{eff}} \!=\! \eta/(1 + \eta\mu)$. So $\mu$ acts as
a \emph{step-size shrinker}, but only when the local step would have
gone large; small steps are unaffected. The expected drift norm should
\emph{shrink} with $\mu$ at the small clients (where $\nabla F_k$ is
noisy and large) more than at the large clients (where the gradient is
already well-estimated).

\textbf{Q5. Observation.}
\auprc{} climbs monotonically: $\mu\!=\!0$ gives $0.636$, $\mu\!=\!0.001$
gives $0.651$, $\mu\!=\!0.01$ gives $0.660$, $\mu\!=\!0.1$ gives
$0.665$. Worst-client recall climbs even faster:
$0.19 \!\to\! 0.41 \!\to\! 0.49 \!\to\! 0.50$. Per-seed std drops by
$\sim$$3\!\times$ between $\mu\!=\!0$ and $\mu\!=\!0.1$.

\textbf{Q6. Reconciliation.}
The drift-norm prediction is verified directly:
\texttt{drift\_norms\_summary.csv} shows mean drift norm at
$\mu\!=\!0$: $0.0145$; at $\mu\!=\!0.1$: $0.0021$ -- a
$7\!\times$ contraction, with the largest contraction at the small
Neuro clients (\texttt{drift\_norms\_per\_client.csv}: ratio
$0.0030/0.0009 \!\approx\! 3.3$ at MICU vs.\
$0.0188/0.0028 \!\approx\! 6.7$ at Neuro Intermediate).

\textbf{Q7. Failure mode.}
At $\mu\!=\!0.1$ the accuracy mean drops slightly ($0.842$ vs.\
$0.878$). The proximal anchor is biasing the per-round model toward
the size-weighted average even when the average is wrong for some
client. If $\mu \!\to\! \infty$ the model freezes at $w_0$; we would
expect a U-shape in any quality metric as $\mu$ grows beyond $0.1$.

\textbf{Q8. What we changed to fix it.}
Three things. First, the \texttt{named\_parameters()} fix above
(otherwise BatchNorm running statistics enter the proximal term and
contaminate $\mu$). Second, we evaluate at multiple $\mu$ values rather
than picking a single one. Third, we report per-seed std so the gain
is read as a confidence statement, not a point estimate.

\textbf{Q9. Deployment guidance.}
$\mu \!=\! 0.1$ is the headline recommendation: best \auroc{}, best
\auprc{}, best worst-client recall, lowest per-seed std. Communication
cost is unchanged from \fedavg{}.

\textbf{Q10. What we would do differently.}
Sweep $\mu$ on a log grid into $\{0.3, 1.0\}$ to find the U-shape
turning point; explore an adaptive per-client $\mu_k$ proportional to
$1/|\mathcal{D}_k|$ so small clients get a stronger anchor.

\subsection{Dossier 3 -- \cvar{}-style aggregation reweighting}
\label{sec:dossier-cvar}

\textbf{Q1. Formal definition.}
After local training in round $t$, sort sampled clients $S_t$ by their
local loss $\hat F_k$ in descending order. Define
$S_t^{\alpha} = \{k \in S_t : \hat F_k \geq Q_{\alpha}(\hat F)\}$,
the upper-$\alpha$ tail. Aggregate uniformly over the tail:
$w_{t+1} = \frac{1}{|S_t^{\alpha}|}\sum_{k \in S_t^{\alpha}} w_{t,k}^{(E)}$.

\textbf{Q2. Mechanism hypothesis.}
\cvar{} aggregation puts all weight on the worst clients, so the
update direction is the gradient of the high-loss tail. Increasing
$\alpha$ should monotonically improve worst-client recall by making
the global step explicitly minimise the tail.

\textbf{Q3. Code location.} \texttt{flopt/aggregators.py},
function \texttt{cvar\_aggregator}. The tail size is computed as
$\lfloor (1-\alpha)\!\cdot\!|S_t| \rfloor$ with a fallback to
``keep the single worst'' when the floor is zero.

\textbf{Q4. Math prediction.}
The expected size of the tail at $|S_t|\!=\!C\!=\!5$ is
$\lfloor 5(1-\alpha) \rfloor$: $\alpha\!=\!0\!\to\!5$,
$\alpha\!=\!0.5\!\to\!2$, $\alpha\!=\!0.75\!\to\!1$,
$\alpha\!=\!0.9\!\to\!0$ (fallback to $1$),
$\alpha\!=\!0.95\!\to\!0$ (fallback to $1$). So we predict
\emph{three distinct regimes}: full average ($\alpha\!=\!0$),
partial tail ($\alpha\!=\!0.5$), single-worst ($\alpha\!\geq\!0.75$),
with $\alpha \!=\! 0.75, 0.9, 0.95$ producing identical updates.

\textbf{Q5. Observation.}
$\alpha\!=\!0.75, 0.9, 0.95$ produce identical \auprc{} ($0.615$),
identical accuracy ($0.866$), identical worst-client recall ($0.24$),
and identical communication ($579$ Mb). Figure~\ref{fig:cvar-alpha-collapse}
shows the collapse explicitly: three points on top of each other.

\textbf{Q6. Reconciliation.}
Exactly the discrete prediction. The ``CVaR knob'' is a
\emph{three-position dial} at $C\!=\!5$, not a continuous one. To get
true CVaR, $C$ must grow.

\textbf{Q7. Failure mode.}
Two failure modes. (i) Discreteness: at $C\!=\!5$, $\alpha$ above
$0.75$ does nothing. (ii) High variance: \cvar{}-$0$ has
$\worstrec{}\!=\!0.26 \!\pm\! 0.21$, meaning $\sim$$40\%$ of seeds
returned $0$. Aggregating the top-loss client at every step lets a
single noisy client steer the global model.

\textbf{Q8. What we changed to fix it.}
We did not change \cvar{} -- we instead reported it as a study object,
so a future paper that wants real CVaR knows to pick $C\!=\!20$ or
larger. The substantive fairness gain came from \fedprox{}-$0.1$, not
\cvar{}.

\textbf{Q9. Deployment guidance.}
Do not use this implementation of \cvar{} in production unless $C \geq
20$. Even then, expect high seed variance and consider a continuous
risk-averse loss (e.g.\ tail averaging on \emph{loss}, not on weights).

\textbf{Q10. What we would do differently.}
Replace the discrete tail with a soft-CVaR weighting
$p_k \propto \mathrm{softmax}(\hat F_k / \tau)$ with $\tau \!\to\! 0$
recovering the hard tail. This makes $\alpha$ continuous, removes the
$C\!=\!5$ collapse, and is differentiable end-to-end.

\subsection{Dossier 4 -- Centralized baseline}
\label{sec:dossier-centralized}

\textbf{Q1. Formal definition.}
Pool all training data $\mathcal{D} = \bigcup_k \mathcal{D}_k$ into a
single training set, train the same TabularMLP for $50$ epochs with
Adam ($\eta\!=\!10^{-3}$, $B\!=\!256$), evaluate per client.

\textbf{Q2. Mechanism hypothesis.}
Centralized is the per-client \emph{ceiling}: any federated method
that does not exceed centralized has either failed to extract the
heterogeneity dividend or has bought it at communication cost. Equally,
centralized's failures (e.g.\ low \auprc{}) are
\emph{intrinsic} to the architecture, not federation.

\textbf{Q3. Code location.}
\texttt{flopt/baselines.py}, function \texttt{centralized\_train}.

\textbf{Q4. Math prediction.}
Centralized minimises the size-weighted population loss
$\sum_k p_k F_k$. This is the same objective as \fedavg{} \emph{at
infinity rounds}, but with no aggregation noise and no client drift.
We therefore predict centralized accuracy $>$ \fedavg{}, with
\auprc{} \emph{not necessarily} better because population-weighting
under-represents the rare class.

\textbf{Q5. Observation.}
Centralized: accuracy $0.908 \!\pm\! 0.004$, \auroc{}
$0.878 \!\pm\! 0.007$, \auprc{} $0.585 \!\pm\! 0.008$,
$\worstrec{}\!=\!0.16 \!\pm\! 0.13$. \fedprox{}-$0.1$ beats centralized
on \auroc{} ($0.913$ vs.\ $0.878$), \auprc{} ($0.665$ vs.\ $0.585$),
and worst-client recall ($0.50$ vs.\ $0.16$).

\textbf{Q6. Reconciliation.}
The accuracy ranking is as predicted (centralized $>$ federated). The
\auprc{} ranking is reversed in our favour: \fedprox{}-$0.1$
\emph{beats} centralized. The mechanism is that centralized is
solving the \emph{wrong} objective for \auprc{} (it's minimising
log-loss on the population mix, where the rare class is sparse), and
\fedprox{}-$0.1$ inadvertently up-weights the rare class because it
prevents the model from over-fitting to the dominant ICU pattern.

\textbf{Q7. Failure mode.}
Centralized's failure mode is the rare class. With base mortality
$\sim$$11\%$ overall and as low as $\sim$$2\%$ at Neuro Intermediate,
log-loss minimisation under-weights deaths and over-confidently
predicts survival. The result is high accuracy, low \auprc{}.

\textbf{Q8. What we changed to fix it.}
We added the proposal-alignment logistic backbone, where we tested
whether a convex objective + careful early stopping could close the
\auprc{} gap. It did not; the issue is class imbalance, not
optimisation pathology.

\textbf{Q9. Deployment guidance.}
Centralized is a comparison object only; pooling clinical data across
ICUs is not a permissible deployment in the proposal's scope.

\textbf{Q10. What we would do differently.}
Re-train centralized with class-balanced loss (focal loss,
class weights) to give it the fairest possible shot at \auprc{}; even
then we expect it to lose to \fedprox{}-$0.1$ on worst-client recall.

\subsection{Dossier 5 -- Local-only baseline}
\label{sec:dossier-local}

\textbf{Q1. Formal definition.}
Train a separate TabularMLP per client on $\mathcal{D}_k$ with
identical hyperparameters, evaluate on $\mathcal{D}_k$'s test fold.
Report per-client metrics and the macro-averaged worst case.

\textbf{Q2. Mechanism hypothesis.}
Local-only is the per-client \emph{floor}: it ignores knowledge
transfer between clients. We expect (a) good performance on big
clients (training set is large enough), (b) catastrophic performance
on small clients (insufficient data), and (c) zero worst-client recall
because at least one tiny client will fail to learn the rare class.

\textbf{Q3. Code location.}
\texttt{flopt/baselines.py}, function \texttt{local\_only\_summary}.

\textbf{Q4. Math prediction.}
A model trained on $\sim$$500$ rows with $10$ deaths is statistically
under-sampled for the rare class: SE on death rate is
$\sqrt{0.02 \cdot 0.98 / 500} \!\approx\! 0.006$, so a 95\% CI on
mortality is $[0.008, 0.032]$ -- consistent with both ``no
mortality'' and the observed $0.02$. The local model has no incentive
to predict deaths; expected sensitivity at small clients is
near-zero.

\textbf{Q5. Observation.}
Aggregate accuracy $0.904$, \auroc{} $0.869$, \auprc{} $0.459$,
$\worstrec{}\!=\!0.0$. The worst-client recall \emph{exactly} matches
the prediction (zero on at least one small client).

\textbf{Q6. Reconciliation.}
The aggregate accuracy ranks above \fedavg{}'s headline because each
local model perfectly fits its own data; the macro-worst recall is
\emph{zero} because at least one local Neuro model never predicts
the rare class. So local-only wins headline accuracy and loses
fairness completely; this is the \emph{strongest} possible argument
for federation.

\textbf{Q7. Failure mode.}
Sample inefficiency on the rare class at small clients. No knowledge
transfer.

\textbf{Q8. What we changed to fix it.}
Federation -- specifically \fedprox{}-$0.1$ -- is exactly the fix.

\textbf{Q9. Deployment guidance.}
Local-only is a comparison object; deployment of nine separate
models is also operationally bad (governance, monitoring, drift).

\textbf{Q10. What we would do differently.}
Allow the local-only baseline to share an unsupervised pretraining
stage; this would give it a better embedding floor and probably a
better worst-client recall.

\subsection{Dossier 6 -- Grid search}
\label{sec:dossier-grid}

\textbf{Q1. Formal definition.}
Discrete Cartesian sweep over $E \!\in\! \{1,2,3\}$,
$C \!\in\! \{5,10,15\}$, $\eta \!\in\! \{0.003, 0.005, 0.01\}$ ($27$
combinations) for the proposal-alignment run; the original run swept
the same axes at lower resolution ($12$ combinations).

\textbf{Q2. Mechanism hypothesis.}
Grid search is the cheapest, most reproducible search. It finds the
best mesh point. If the loss surface is smooth and the optimum is
near a mesh point, grid is competitive; if the optimum is between
mesh points or in a non-axis-aligned valley, grid will lose to
stochastic search.

\textbf{Q3. Code location.}
\texttt{flopt/search.py}, function \texttt{grid\_search}; results in
\texttt{outputs/full\_mimic\_iv\_proposal\_alignment/search/proposal\_grid\_search.csv}.

\textbf{Q4. Math prediction.}
With $27$ points and $\eta$-axis only at three levels
$\{0.003, 0.005, 0.01\}$, the closest mesh point to the DE optimum
$\eta\!\approx\!0.0061$ is $\eta\!=\!0.005$ (offset $20\%$). The
expected gap between grid and DE is roughly the slope of loss at
$\eta\!=\!0.005$ times $0.0011$. From the landscape figure, the loss
slope at the optimum is small, so we predict the gap is real but
modest.

\textbf{Q5. Observation.}
Grid best: \auprc{} $0.598$, accuracy $0.901$, comm $359$ Mb. GA best:
\auprc{} $0.629$, accuracy $0.867$, comm $318$ Mb. Grid won on
accuracy (mesh point hits a local max), GA won on \auprc{} and on
communication.

\textbf{Q6. Reconciliation.}
The off-mesh $\eta$ found by DE is the precise mechanism we identify;
grid is exactly $\eta\!=\!0.005$ at its best, DE is at
$\eta\!=\!0.00614$ -- a 23\% offset that the grid cannot represent
without a 4-axis grid (108 combinations).

\textbf{Q7. Failure mode.}
Quantisation. Any fine optimum that does not align with the mesh is
unreachable.

\textbf{Q8. What we changed to fix it.}
Added DE search alongside grid; reported both and let the LP duality
analysis pick the operating point.

\textbf{Q9. Deployment guidance.}
Use grid search for sanity-checking the order of magnitude of each
hyperparameter. Use DE/GA for refinement. The combined cost
($27 + 90$ evaluations) is well within a single GPU-day on \mimic{}.

\textbf{Q10. What we would do differently.}
Adopt a halving (Successive Halving / Hyperband) strategy that
allocates more rounds to promising candidates, so the search budget
goes to the right place automatically.

\subsection{Dossier 7 -- Differential Evolution / Genetic Algorithm}
\label{sec:dossier-de}

\textbf{Q1. Formal definition.}
\texttt{scipy.optimize.differential\_evolution} on a $3$-D box
($E$, $C$, $\eta$) with \texttt{popsize}$=15$, \texttt{maxiter}$=15$,
fitness $= \text{loss} + \beta_C \cdot (\text{comm}/10^9)$,
$\beta_C \!=\! 1$. Each evaluation is one $80$-round
\fedavg{} run on $30$ evaluation clients.

\textbf{Q2. Mechanism hypothesis.}
DE is a stochastic search that can find off-grid optima. We expect it
to find an $\eta$ that is not in the grid set, and we expect it to
\emph{drive communication down} because the fitness function
explicitly penalises bytes.

\textbf{Q3. Code location.} \texttt{flopt/search.py}, function
\texttt{ga\_search}; convergence trace in
\texttt{outputs/full\_mimic\_iv\_proposal\_alignment/search/proposal\_ga\_best\_so\_far.csv}.

\textbf{Q4. Math prediction.}
DE explores $90$ candidates ($15 \times 6$) over the same box as grid.
Random sampling at this density gives $\sim$$0.5\%$ probability that
a uniform candidate lands within $1\%$ of the grid optimum, but DE's
adaptive recombination means the probability is far higher; we
predict DE will visit the local minimum and converge by iteration
$\sim$$10$.

\textbf{Q5. Observation.}
Best fitness reached at iteration $9$ ($31$st evaluation in the
\texttt{best\_so\_far} curve); plateau thereafter. Best operating
point: $E\!=\!1$, $C\!=\!5$, $\eta\!=\!0.00614$ -- exactly off-mesh.
\auprc{} $0.629$, communication $318$ Mb (vs.\ grid's $359$).

\textbf{Q6. Reconciliation.}
The convergence curve in Figure~\ref{fig:ga-convergence} shows the
exact predicted shape: rapid early decrease, plateau around iteration
$10$. The off-mesh $\eta$ is the precise mechanism for the
grid-vs-DE \auprc{} gap.

\textbf{Q7. Failure mode.}
Stochasticity: DE depends on random initialisation. With a different
seed, the operating point could shift. We did not seed DE; the next
run could land elsewhere on the loss-vs-comm Pareto.

\textbf{Q8. What we changed to fix it.}
Reported the search history in full so a re-run can be compared
directly.

\textbf{Q9. Deployment guidance.}
For \mimic{}, DE finds a Pareto-better operating point than grid
(loss \emph{and} bytes lower) within $\sim$$30$ evaluations. Run DE
and validate the operating point on a held-out \emph{client}
(not just held-out rows).

\textbf{Q10. What we would do differently.}
Replace DE with a multi-objective search (NSGA-II or BoTorch's
qEHVI) that returns the entire loss-vs-comm Pareto, not a single
weighted compromise. Seed the search for reproducibility.

\subsection{Dossier 8 -- LP duality (the communication shadow price)}
\label{sec:dossier-lp}

\textbf{Q1. Formal definition.}
Let $\mathcal{X} \!=\! \{x_i\}_{i=1}^{N}$ be the candidate set
(GA history, $N\!\in\!\{48,90\}$). Each $x_i$ has loss $\ell_i$ and
communication $b_i$. For budget $B$, solve
\[
\min_{p \in \Delta^{N}} \sum_i p_i \ell_i \quad \text{s.t.}
\quad \sum_i p_i b_i \leq B.
\]
The dual variable on the budget constraint is the \emph{shadow price}
$\lambda(B) \!=\! -\partial \ell^* / \partial B$.

\textbf{Q2. Mechanism hypothesis.}
$\lambda(B)$ is the marginal value of one byte. When the budget binds
($B$ is small enough that any cheaper candidate is preferred even at
slightly worse loss), $\lambda > 0$. When the budget is slack, $\lambda
\!=\! 0$ and any extra byte buys nothing.

\textbf{Q3. Code location.}
\texttt{flopt/duality.py}; results in
\texttt{outputs/full\_mimic\_iv\_proposal\_alignment/lp/} (twelve
budgets dense, twelve sparse, plus
\texttt{kkt\_diagnostics.csv}).

\textbf{Q4. Math prediction.}
The LP is a $N$-simplex with one inequality constraint. The optimal
$p$ is concentrated on at most two vertices (the two
loss-vs-comm-Pareto-adjacent points around the budget). $\lambda$ is
the slope between them. Above the largest $b_i$, $\lambda \!=\! 0$;
below the smallest $b_i$, the LP is infeasible.

\textbf{Q5. Observation.}
Dense (MLP) LP: $\lambda$ varies over thirteen orders of magnitude --
$\sim$$10^{-7}$ at the active region, $\sim$$10^{-20}$ at the slack
region. The transition is the largest-budget GA candidate. Sparse
(logistic) LP: same shape, shifted left by two orders of magnitude
(the logistic candidates communicate $\sim$$50\!\times$ less).

\textbf{Q6. Reconciliation.}
Exactly the predicted two-region shape. KKT diagnostics
(\texttt{kkt\_diagnostics.csv}) confirm the LP is solved cleanly --
all $24$ programs report \texttt{primal\_feasible=True},
\texttt{dual\_feasible=True}, complementary slackness $\leq 8\!\times\!
10^{-10}$.

\textbf{Q7. Failure mode.}
Two failure modes. (i) The LP only ranks operating points it has seen;
if the true Pareto frontier is not in $\mathcal{X}$, the LP cannot
recommend it. (ii) The shadow price is a slope, not a step; small
budget changes around the active region have linear effect, but the
\emph{operating recommendation} is a discrete vertex flip.

\textbf{Q8. What we changed to fix it.}
Reported the LP at twelve budgets so the entire shadow-price curve
is visible. Reported KKT residuals so the LP solution is auditable.

\textbf{Q9. Deployment guidance.}
$\lambda(B)$ at the realistic budget for \mimic{} ($\sim$$300$ Mb) is
$\sim$$10^{-9}$, meaning a million extra bytes buys $\sim$$10^{-3}$
loss reduction. That is operationally negligible. The conclusion
``stop spending bytes'' is the LP's own recommendation.

\textbf{Q10. What we would do differently.}
Run the LP with the GA candidates \emph{and} the centralized point
\emph{and} the local-only point so the LP can choose the entire
spectrum, not just federated configurations.

\subsection{Dossier 9 -- Sparsity / top-$k$ communication}
\label{sec:dossier-sparsity}

\textbf{Q1. Formal definition.}
At each round, after local training, keep only the largest $k\%$ of
update entries by absolute value, zero the rest, communicate only the
survivors with their indices. Reconstruct the dense update at the
server.

\textbf{Q2. Mechanism hypothesis.}
Two distinct mechanisms. (a) Natural sparsity: if updates are already
near-sparse, top-$k$ is a free compression. (b) Lossy compression: if
updates are dense, top-$k$ is information loss bought for budget.

\textbf{Q3. Code location.}
\texttt{flopt/sparsity.py}; results in
\texttt{outputs/full\_mimic\_iv\_proposal\_alignment/lp/sparsity\_summary.csv}
and \texttt{dense\_vs\_sparse\_shadow\_price.csv}.

\textbf{Q4. Math prediction.}
For Adam-trained MLPs we expect dense updates ($95\%$+ nonzero); the
top-$k$ at $k\!=\!1$ retains only the largest $1\%$. The loss penalty
should be roughly the variance of the discarded entries times their
direction; for a moderately-trained MLP this is dominant in early
rounds and shrinks as the model approaches convergence.

\textbf{Q5. Observation.}
Natural nonzero rate $0.958$--$0.977$ (i.e.\ \emph{no} natural
sparsity). Top-$1\%$ achieves $\sim$$49\!\times$ compression but at
the price of throwing away $99\%$ of the gradient direction. The
Pareto frontier favours the dense MLP at the operating budgets.

\textbf{Q6. Reconciliation.}
Exactly the lossy-compression mechanism. There is no free
compression dividend on \mimic{}.

\textbf{Q7. Failure mode.}
At $k\!=\!1\%$ the model is essentially doing $1\%$-of-coordinates
SGD; convergence is slow (more rounds needed) and unstable.

\textbf{Q8. What we changed to fix it.}
Reported sparsity as a study object (LP duality), not as a deployment
recommendation.

\textbf{Q9. Deployment guidance.}
On \mimic{}, do not deploy top-$1\%$. The LP at the realistic budget
prefers logistic backbone over compressed MLP. If you need
compression, use error feedback (e.g.\ EF-SGD) which corrects the
discarded entries on the next round.

\textbf{Q10. What we would do differently.}
Add error feedback; sweep $k \!\in\! \{0.1\%, 0.5\%, 5\%\}$; compare
to randomised projections (sketching).

\subsection{Dossier 10 -- Dirichlet stress test}
\label{sec:dossier-dirichlet}

\textbf{Q1. Formal definition.}
Construct $K\!=\!30$ synthetic clients by drawing label proportions
from $\mathrm{Dir}(\beta\,\mathbf{1}_2)$ and sampling
$\sim$$|\mathcal{D}|/K$ rows accordingly. Sweep
$\beta \!\in\! \{0.1, 0.5, 1.0, \infty\}$.

\textbf{Q2. Mechanism hypothesis.}
$\beta$ controls heterogeneity. $\beta\!\to\!\infty$ recovers IID
($\mathrm{Dir}(\infty)$ is a point mass at the uniform); $\beta\!\to\!0$
gives clients with vastly different label rates. We predict that
\fedavg{} degrades smoothly with falling $\beta$, \fedprox{} is more
robust, and \cvar{} (with its small effective $\alpha$) is in
between.

\textbf{Q3. Code location.}
\texttt{flopt/dirichlet.py}, results in
\texttt{outputs/full\_mimic\_iv\_proposal\_alignment/metrics/dirichlet\_beta\_summary.csv}
(method $\times$ $\beta$ $\times$ seed).

\textbf{Q4. Math prediction.}
\fedavg{} accuracy at $\beta\!=\!0.1$ should drop noticeably below
the $\beta\!=\!1$ point because clients have very skewed label rates
and the size-weighted average is no longer close to centralized.
\fedprox{}'s anchor should buffer the drop.

\textbf{Q5. Observation.}
At $\beta\!=\!0.1$: \fedavg{} acc $0.870 \!\pm\! 0.044$, \fedprox{}
$0.870 \!\pm\! 0.045$. At $\beta\!=\!1.0$: \fedavg{} $0.876$,
\fedprox{} $0.876$. Per-seed std is high ($0.04$--$0.11$) at
$\beta\!=\!0.1$, low ($0.01$) at $\beta\!=\!\infty$.

\textbf{Q6. Reconciliation.}
Mean accuracy is robust. The stress shows up in \emph{seed variance},
not mean. The mechanism is that \mimic{}'s feature space is
discriminative enough that extreme label imbalance is still learnable
in the mean -- but particular seeds (where the synthetic clients
land in unfortunate corners of feature space) fail completely. This
is exactly what real-world heterogeneity looks like: the mean is
fine, the tail is brutal.

\textbf{Q7. Failure mode.}
Per-seed instability at low $\beta$.

\textbf{Q8. What we changed to fix it.}
Reported per-seed std prominently. \fedprox{}-$0.1$ has the lowest
across-seed std at every $\beta$.

\textbf{Q9. Deployment guidance.}
On a heterogeneous deployment, demand per-seed std as part of the
SLA, not just mean.

\textbf{Q10. What we would do differently.}
Run more seeds at $\beta\!=\!0.1$ ($30$+ rather than $10$) to
characterise the tail of failure cases.

\subsection{Dossier 11 -- Loss-landscape interpolation}
\label{sec:dossier-landscape}

\textbf{Q1. Formal definition.}
Following Li et al.\ (2018), pick $w_{\text{init}}, w_{\text{best}},
w_{\text{final}}$. Plot $L(\alpha) \!=\! \text{loss}((1\!-\!\alpha)
w_{\text{init}} \!+\! \alpha w_{\text{final}})$ for
$\alpha \!\in\! [-0.5, 1.5]$ (1D), and
$L(\alpha_1, \alpha_2)$ in two random filter-normalised directions on
a $25\!\times\!25$ grid (2D).

\textbf{Q2. Mechanism hypothesis.}
A smooth, convex-looking landscape implies that grid search finds the
optimum reliably. A rough, non-convex landscape implies that
stochastic search wins because grid mesh points may not align with
optima. The MLP landscape is a classical neural-net surface; we
expect roughness in both backbones, but \emph{less} for logistic
(convex objective).

\textbf{Q3. Code location.}
\texttt{flopt/landscape.py}; data in
\texttt{outputs/full\_mimic\_iv\_proposal\_alignment/landscape/}.

\textbf{Q4. Math prediction.}
Logistic regression is convex, so its 1D and 2D slices should be
strictly convex. The MLP is non-convex; we predict 1D shows a smooth
descent with no obvious bumps, but 2D may show non-axis-aligned
curvature.

\textbf{Q5. Observation.}
Figure~\ref{fig:landscape-1d} (1D, both backbones): smooth bowl
in both. MLP minimum is at $\alpha\!=\!0.85$; logistic at $\alpha\!=\!1$
exactly. Figure~\ref{fig:landscape-2d} (2D): logistic is a
bowl; MLP shows a slight ridge in one direction.

\textbf{Q6. Reconciliation.}
The 1D MLP minimum at $\alpha\!=\!0.85$ (not $1$) is the precise
mechanism for the grid-vs-DE gap: the best round of \fedavg{} is
closer to $0.85$ along the init-to-final axis, so the global model
at the end of training has \emph{overshot}. DE finds the
``earlier-stop'' operating point because its fitness function
penalises bytes (and bytes correlate with rounds, hence with $\alpha$).

\textbf{Q7. Failure mode.}
The two random filter-normalised directions are not the
loss-axis-aligned principal directions; the 2D figure may
under-represent the curvature. Hessian-based directions would be
truer.

\textbf{Q8. What we changed to fix it.}
Reported both backbones (logistic for sanity; MLP for the actual
deployment case) so the cross-check is in the same figure.

\textbf{Q9. Deployment guidance.}
The 1D landscape minimum at $\alpha\!=\!0.85$ is a deployment
recommendation: terminate training $\sim$$15\%$ before the final
round, even if validation loss is still decreasing.

\textbf{Q10. What we would do differently.}
Use Hessian eigenvector directions (Lanczos) in the 2D figure;
report Hessian top-$k$ eigenvalues at the best checkpoint.

\subsection{Dossier 12 -- Logistic regression backbone (the convex sanity)}
\label{sec:dossier-logistic}

\textbf{Q1. Formal definition.}
A single linear layer on the $30$-dimensional feature vector
($\sigma(w^\top x + b)$). Same federated training pipeline as the
MLP -- only the model class changes.

\textbf{Q2. Mechanism hypothesis.}
Logistic is a convex objective, $50\!\times$ smaller in parameters,
and produces a calibrated probability if trained with cross-entropy.
We expect (a) much lower communication, (b) somewhat lower \auprc{}
(less expressive), (c) better calibration (single-layer, less
overfitting), (d) lower per-seed variance.

\textbf{Q3. Code location.}
\texttt{flopt/models.py}, class \texttt{LogisticModel}; results in
\texttt{outputs/full\_mimic\_iv\_proposal\_alignment/metrics/logreg\_method\_summary.csv}.

\textbf{Q4. Math prediction.}
Logistic communication is $\sim$$30 \cdot 4 \cdot N_{\text{rounds}}$
bytes per client per round (plus bias and gradient stats). MLP
communication is $\sim$$30 \cdot 64 + 64 \cdot 64 + 64 \cdot 2 \approx
6{,}400$ parameters $\!\cdot\!$ rounds, i.e.\ $\sim$$50\!\times$ more.
We predict logistic at $\sim$$6$ Mb total comm, MLP at $\sim$$300$
Mb -- a factor of $50$.

\textbf{Q5. Observation.}
Logistic \fedavg{}: \auprc{} $0.601$, comm $6.4$ Mb, $\worstrec{}\!=\!
0.45$. Logistic \fedprox{}-$0.1$: \auprc{} $0.612$, comm $6.4$ Mb,
$\worstrec{}\!=\!0.45$. MLP \fedprox{}-$0.1$: \auprc{} $0.665$,
comm $428$ Mb. Worst-client recall is \emph{nearly the same} as MLP
($0.45$ vs.\ $0.50$). Per-seed std on logistic is $0.001$--$0.010$
(extremely low).

\textbf{Q6. Reconciliation.}
The communication ratio is $428/6.4 \!\approx\! 67$ -- close to the
predicted $50\!\times$. The \auprc{} drop is $0.053$, modest. The
worst-recall is robust -- the proximal anchor's effect is the same
mechanism in both backbones. Logistic's per-seed std is exceptional:
the convex surface has a unique optimum, and federation reaches it
reliably.

\textbf{Q7. Failure mode.}
\auprc{} ceiling: logistic cannot capture interaction effects between
features (e.g.\ ``high heart rate \emph{and} low BP \emph{and} elevated
lactate''). On a richer feature set (raw waveforms, free text), the
gap to MLP would widen.

\textbf{Q8. What we changed to fix it.}
Reported both backbones in every chart, so the $\Delta\auprc{}\!=\!
0.053$ at $50\!\times$ less communication is the reader's
operational decision.

\textbf{Q9. Deployment guidance.}
For a bandwidth-constrained edge ICU (e.g.\ rural hospital), logistic
\fedprox{}-$0.1$ is the deployment recommendation. For a
bandwidth-rich tertiary centre, MLP \fedprox{}-$0.1$ is the
recommendation.

\textbf{Q10. What we would do differently.}
Try a $2$-layer ($30 \!\to\! 16 \!\to\! 1$) intermediate model and
sweep the depth-vs-bandwidth Pareto.

\paragraph{Dossier summary.}
Twelve algorithms, twelve ten-question dossiers, one consistent
finding: \fedprox{}-$0.1$ is the substantive winner, the LP says the
budget is in the slack region (so dollars and bytes should be spent
elsewhere), the landscape says terminate $\sim$$15\%$ early, and the
logistic backbone is a credible $50\!\times$-cheaper alternative for
bandwidth-constrained deployment. Everything else is a study object.
